%====================================================================
% Universidade do Estado do Pará
% Centro de Ciências Sociais e Educação - CCSE
% Curso: Licenciatura em Física
% Modelo para elaboração de monografia referentes aos trabalhos de
% conclusão de curso de Licenciatura em Física
% Autores(as): 
% Felipe Vera da Rocha, graduando do curso de Licenciatura em Física do ano de 2019 da Universidade do Estado do Pará (UEPA), Campus XX - Castanhal. 
% Athirson Fernando Borges Lima, graduando do curso de Licenciatura em Física do ano de 2019 da Universidade do Estado do Pará (UEPA), Campus XX - Castanhal.
% Professor Dr. Erick Elisson Hosana Ribeiro, coordenador de TCC do curso de Licenciatura em Física da Universidade do Estado do Pará (UEPA).
% Professor Dr. Alessandre Sampaio da Silva, coordenador do curso de Licenciatura em Física da Universidade do Estado do Pará (UEPA).
%====================================================================
%                        Tipo do Documento 
%====================================================================

\documentclass[12pt,				% tamanho da fonte
	openright,			% capítulos começam em página ímpar (insere página vazia caso preciso)
	oneside,			% para impressão em frente e verso. Oposto a oneside
	a4paper,			% tamanho do papel. 
	% -- opções da classe abntex2 --
	chapter=TITLE,		% títulos de capítulos convertidos em letras maiúsculas
	%section=TITLE,		% títulos de seções convertidos em letras maiúsculas
	%subsection=TITLE,	% títulos de subseções convertidos em letras maiúsculas
	%subsubsection=TITLE,% títulos de subsubseções convertidos em letras maiúsculas
	% -- opções do pacote babel --
	sumario=abnt-6027-2012, % Sumario de acordo com as normaisda ABNT
	english,			% idioma adicional para hifenização
	spanish,			% idioma adicional para hifenização
	brazil				% o último idioma é o principal do documento
]{abntex2}

%====================================================================
% O comando \usepackage{estilo} tem por finalidade carregar os pacotes do LaTeX bem como definir as configurações do texto da monografia. 
%====================================================================

\usepackage{estilo}

%====================================================================
%                   Informações Gerais da Monografia
%====================================================================
%====================================================================
% Este documento foi montado de tal forma que as informações gerais
% vão ser organizadas automatimanete pelo texto. Para isso preencha
% corretamente os dados abaixo. 
%====================================================================

% INFORMAÇÕES DA UNIVERSIDADE E CURSO
\universidade{Universidade do Estado do Pará - UEPA} % Não precisa mudar
\centro{Centro de Ciências Sociais e Educação - CCSE} % Não precisa mudar
\curso{Curso de Licenciatura em Física} 
\formacao{Licenciado(a) em Física} 

% INFORMAÇÕES DA MONOGRAFIA
\titulo{Colocar Título}
\autor{Como nome do autor(a) ou  autores(as)}
\tipotrabalho{Trabalho de Conclusão de Curso} % Não precisa mudar
\local{Cidade - PA} % Ex: Castanhal - PA
\data{Ano}
\datadadefesa{Data da Defesa: xx de xxxxx de xxxxx}

\preambulo{Trabalho de Conclusão de Curso apresentado como requisito parcial para obtenção do grau de \imprimirformacao, da \imprimiruniversidade. \\ 

Orientador: \imprimirorientador.} % Não precisa mudar

% INFORMAÇÕES PARA FICHA CATALOGRÁFICA
\nomedoprimeiroautor{Nome Completo}  
%\nomedosegundoautor{Nome Completo} 

% caso seu trabalho possua mais de um autor apenas preencha os dado e
% retire os comentários dos campos anteriores e na fica
% catalográfica

\palavraschave{%
1. Palavra-chave1.
2. Palavra-chave2.
3. Palavra-chave3.
}

% INFORMAÇÕES DO ORIENTADOR E DO COORIENTADOR 
\orientador{Prof. Dr. Nome Completo} 
\faculdadedoorientador{Informar}
\coorientador{Prof. Dr. Nome Completo} % Não Obrigatório
\faculdadedocoorientador{Informar} % Não Obrigatório

% Caso o seu trabalho possua um coorientador é necessário habilitar
% o seu campos de informações na folha de aprovação.

% INFORMAÇÕES DOS MEMBROS DA BANCA
\primeiromembrodabanca{Prof. Dr. Nome Completo}
\faculdadedoprimeiromembrodabanca{Informar}
\segundomembrodabanca{Prof. Dr. Nome Completo}
\faculdadedosegundomembrodabanca{Informar}
\terceiromembrodabanca{Prof. Dr. Nome Completo}
\faculdadedoterceiromembrodabanca{Informar}

%====================================================================
%                   Documento Principal - Início
%====================================================================

\begin{document}
\selectlanguage{brazil}
\pretextual

%====================================================================
% CAPA 
%====================================================================

\thispagestyle{empty}
\input{1 - Elementos Pré-Textuais/1.1 - Capa}

%====================================================================
% FOLHA DE ROSTO
%====================================================================

\thispagestyle{empty}
\input{1 - Elementos Pré-Textuais/1.2 - Folha de Rosto}

%====================================================================
% FICHA CATALOGRÁFICA
%====================================================================

\input{1 - Elementos Pré-Textuais/1.3 - Ficha Catalográfica}

% ATENÇÃO !!!
% ESTA PAGINA FOI CRIADA COMO UM ""POSSÍVEL MODELO"" DE FICHA CATALOGRAFICA.
% ENTRETANTO ESTE DOCUMENTO DEVE SER ELABORADO PELA BIBLIOTECA DA INSTITUIÇÃO.
%RECOMENDAMOS QUE BUSQUE MELHORES INFORMAÇÕES PARA A ORGANIZAÇÃO DESTE PARTE DO DOCUMENTO NA BIBLIOTECA.

%====================================================================
% FOLHA DE APROVAÇÃO
%====================================================================

\thispagestyle{empty}
\input{1 - Elementos Pré-Textuais/1.4 - Fola de Aprovação}

%====================================================================
% DEDICATÓRIA
%====================================================================

\thispagestyle{empty}
\input{1 - Elementos Pré-Textuais/1.5 - Dedicatória}

%====================================================================
% AGRADECIMENTOS
%====================================================================

\thispagestyle{empty}
\input{1 - Elementos Pré-Textuais/1.6 - Agradecimentos}

%====================================================================
% EPÍGRAFE
%====================================================================

\thispagestyle{empty}
\input{1 - Elementos Pré-Textuais/1.7 - Epígrafe}

%====================================================================
% RESUMO (Portugês)
%====================================================================

\thispagestyle{empty}
\input{2 - Resumos/2.1 - Resumo em Portugês}

%====================================================================
% RESUMO (Inglês)
%====================================================================

\thispagestyle{empty}
\input{2 - Resumos/2.2 Resumo em Inglês}

%====================================================================
% LISTA DE ILUSTRAÇÕES
%====================================================================

\thispagestyle{empty}
\listoffigures*
\cleardoublepage

%====================================================================
% LISTA DE TABELAS
%====================================================================

\thispagestyle{empty}
\listoftables*
\cleardoublepage

%====================================================================
% LISTA DE ABREVIATURAS E SIGLAS
%====================================================================

\input{1 - Elementos Pré-Textuais/1.8 - Lista de Abreviaturas de Siglas}

%====================================================================
% LISTA DE SIMBOLOS
%====================================================================

\input{1 - Elementos Pré-Textuais/1.9 - Lista de Símbolos}

%====================================================================
% SUMÁRIO
%====================================================================

\thispagestyle{empty}
\tableofcontents*\newpage


%====================================================================
%                     ELEMENTOS TEXTUAIS
%====================================================================

\textual
% Estilo de cabeçalho e rodapé personalizado
\pagestyle{cabe_uepa}
\setcounter{page}{1}


%====================================================================
% CAPÍTULO 01
%====================================================================

\input{3 - Elementos Textuais/3.1 - Capítulo 1}

%====================================================================
% CAPÍTULO 02
%====================================================================

\input{3 - Elementos Textuais/3.2 - Capítulo 2}

%====================================================================
% CAPÍTULO 03
%====================================================================

\input{3 - Elementos Textuais/3.3 - Capítulo 3}

%====================================================================
%                     ELEMENTOS PÓS-TEXTUAIS
%====================================================================

\postextual


%====================================================================
% REFERÊNCIAS BIBLIOGRÁFICAS
%====================================================================

\bibliography{referencias.bib}

%====================================================================
% APÊNDICES
%====================================================================

\input{4 - Elementos Pós - Textuais/4.1 - Apêndices}

%====================================================================
% ANEXOS
%====================================================================

\input{4 - Elementos Pós - Textuais/4.2 - Anexos}

%====================================================================
% FINAL 
%====================================================================

\newpage
\includepdf{Imagens para as Capas/Capa Final do Curso de Física.pdf}


\end{document}






